% IEEE Network magazine article.
% Special Issue: Quantum-AI Synergies Empowering IoT Network Resilience.
% HARD LIMITS: 4500 words (intro->conclusion), <=6 figures+tables combined,
% no equations (<=3 simple ones with GE consent), 15 references.
% House style: no dashes as punctuation (use commas / parentheses / colons).
\documentclass[journal]{IEEEtran}

\usepackage{graphicx}
\usepackage{booktabs}
\usepackage{tikz}
\usetikzlibrary{arrows.meta, positioning, shapes.geometric}
\usepackage[hidelinks]{hyperref}

\begin{document}

\title{Self-Healing IoT Networks: A Quantum-AI Closed-Loop Framework for Resilience}

\author{Do Phuc Hao and collaborators% TODO author block + affiliations (CAIRA / DAU)
\thanks{Manuscript submitted to the IEEE Network Special Issue on Quantum-AI Synergies Empowering IoT Network Resilience.}}

\maketitle

\begin{abstract}
% ~150 words. Problem: IoT networks fragment under cascading faults and targeted attacks.
% Gap: reactive, myopic repair leaves them fragile to the next failure. Idea: a closed loop
% that couples a lightweight AI perception layer (detect and localize degradation) with a
% quantum-inspired optimization layer (recompute a resilient configuration under an energy
% budget). Synergy: the AI risk map warm-starts and prunes the quantum-inspired search, and
% recovered configurations retrain the detector. Evidence: two case studies show the loop
% restores worst-case survivability that myopic repair cannot, and that AI guidance speeds
% healing. Honest: the quantum-inspired optimizer is competitive with strong metaheuristics,
% and its probabilistic substrate is what makes the AI fusion natural.
\textit{TODO abstract.}
\end{abstract}

\begin{IEEEkeywords}
IoT resilience, self-healing networks, quantum-inspired optimization, edge AI, network robustness.
\end{IEEEkeywords}

% ======================================================================
\section{Introduction}
% ~600 words. Resilience as the central theme. Cascading failures + attacks in IoT. Why
% classical reactive healing is insufficient (myopic, leaves single points of failure). The
% quantum-AI synergy thesis. Contributions: (1) closed-loop self-healing framework,
% (2) AI<->QIO synergy mechanism, (3) two case studies, (4) honest portfolio comparison.

% ======================================================================
\section{Why Quantum-AI Synergy for Resilience}
% ~600 words. Background, accessible. Resilience challenges in IoT (topology robustness,
% jamming, energy limits). What AI brings (fast detection/prediction). What quantum-inspired
% optimization brings (global combinatorial search, GPU-batchable, anytime). Why the
% *synergy* (not either alone): the probabilistic Q-bit substrate is the natural carrier for
% AI priors. Table 1 = taxonomy mapping resilience challenges -> quantum-AI mechanisms.

% ======================================================================
\section{The Self-Healing Closed Loop}
% ~900 words. The framework: Sense -> Detect/Predict (AI) -> Reoptimize (QIO) -> Actuate ->
% Learn. Fig 1 = architecture (TikZ). The two synergy channels: AI->QIO (warm-start + search
% pruning via Q-bit amplitudes) and QIO->AI (recovered configs as feedback). Solver-agnostic
% but QIO-native prior fusion.

% ======================================================================
\section{Case Study 1: Reconfiguration under Cascading Failures}
% ~700 words. Setup, worst-case survivability objective (2-edge-connected core + lambda2)
% under an energy budget. Fig 2 = survivable fraction + budget utilisation (greedy under-heals).
% Fig 4 = synergy ablation (AI warm-start convergence curves / healing latency). Honest table.

% ======================================================================
\section{Case Study 2: Reallocation under Jamming Surge}
% ~700 words. Edge IoT, jamming knocks out links + traffic surge. Reroute/reallocate to
% maximize served demand under power. Fig 3 = served demand + tail latency vs baselines.

% ======================================================================
\section{Discussion and Open Challenges}
% ~500 words. What the synergy buys (faster, budget-aware, resilient healing). Honest limits
% (metaheuristics competitive; quantum-inspired not quantum-supreme). Path to real QPU.
% Open challenges: scalability, distributed/federated healing, security of the loop.

% ======================================================================
\section{Conclusion}
% ~150 words.

\begin{thebibliography}{15}
% <=15 refs. TODO.
\bibitem{ref1} TODO.
\end{thebibliography}

\end{document}
